\documentclass[11pt,a4paper]{extarticle}

% Encoding and fonts
\usepackage[T1]{fontenc}
\usepackage[utf8]{inputenc}
\usepackage{lmodern}

% Geometry and layout
\usepackage[margin=2.5cm]{geometry}
\setlength{\parskip}{0.6em}
\setlength{\parindent}{0pt}

% Maths
\usepackage{amsmath,amssymb}

% Tables
\usepackage{booktabs}
\usepackage{array}

% Colours and links
\usepackage[dvipsnames]{xcolor}
\usepackage[
colorlinks=true,
linkcolor=MidnightBlue,
citecolor=MidnightBlue,
urlcolor=MidnightBlue
]{hyperref}

% Better lists
\usepackage{enumitem}
\setlist[itemize]{topsep=0.3em,itemsep=0.2em}
\setlist[enumerate]{topsep=0.3em,itemsep=0.2em}

% Title info
\title{Backtest of ICES Advice for Selected Northeast Atlantic Stocks}
\author{%
	Laurence Kell\\
	\small Sea++
}
\date{\small \today}

\begin{document}
	
	\maketitle
	
	\section*{Summary}
	
	This project will perform a backtest of fisheries advice for i) Celtic Sea cod, ii) Celtic Sea whiting, iii) Celtic Sea haddock, iv) Irish Sea whiting, v) Irish Sea cod, and vi) Northeast Atlantic mackerel.
	
	To conduct the backtest, stock assessments will be rerun annually over a historical period, i.e.\ as a retrospective analysis by peeling back one year at a time. The advice rule that was in operation at the time will then be applied each year, and the stock projected forward under perfect implementation of the resulting catch advice. Outputs will quantify how stocks would have evolved if management had followed scientific advice; these counterfactual outcomes will be compared with realised catches, stock status, and the risks associated with overfishing.
	
	\section*{Aims and Objectives}
	
	The study has four linked objectives:
	
	\begin{enumerate}
		\item \textbf{Reconstruct} the annual perception of stock status for each stock using only the data and assessment methods available at the time.
		\item \textbf{Reproduce} the agreed harvest control rule (e.g. ICES advice rule) in operation in each historical year, including the standard MSY rule for category~1 stocks and any stock-specific management-plan rules where relevant.
		\item \textbf{Simulate} forward from each assessment year for perfect-management, where realised catches match the advised catch.
		\item \textbf{Compare} counterfactual outcomes with realised history to evaluate stock status and rebuilding, yield, stability, and the consequences of deviations from advice.
	\end{enumerate}
	
	\section*{Scope}
	
	The work covers six candidate stocks that are managed based on ICES quantitative assessments (Category~1 stocks) under an annual catch advice framework. The week~1 screening phase may refine this list.
{\footnotesize	
\begin{table}[h]
	\centering
	\small
	\caption{Study stocks}
	\begin{tabular}{@{}llll@{}}
		\toprule
		Stock & ICES area & Assessment & Advice basis \\
		\midrule
		Celtic Sea cod &
		7.e--k &
		SAM age-structured &
		ICES MSY approach \\
		
		Celtic Sea whiting &
		7.b--c, 7.e--k &
		Analytical age-based &
		ICES MSY approach \\
		
		Celtic Sea haddock &
		7.b--k &
		SAM age-structured &
		ICES MSY approach \\
		
		Irish Sea whiting &
		7.a &
		SAM age-structured &
		ICES MSY approach \\
		
		Irish Sea cod &
		7.a &
		Age-structured analytical &
		ICES MSY approach \\
		
		NE Atl.\ mackerel &
		1--8, 14, 9.a &
		Integrated age-structured &
		MSY rule (Coastal States) \\
		\bottomrule
	\end{tabular}
	\end{table}}
		
	\section*{Methods}
	
	\paragraph{Initial screening (week~1)}
	An initial screening exercise will obtain the most recent ICES assessments and historical advice from ICES Stock Assessment Graphs (estimates of stock status, advised catch, TAC, and reported catch). Consistency of scientific advice, advice--management mismatch (i.e.\ whether advice is followed), and the stability of stock assessment estimates provided by ICES will be evaluated. Following this exercise, the final case studies for the backtest may be modified.
	
	\paragraph{Benchmark projections}
	Once the case-study stocks are confirmed, short-term projections will be run from the most recent assessment at target fishing mortalities (e.g.\ \(F_{\mathrm{MSY}}\)) and outcomes compared with the current assessment as a benchmark.
	
	\paragraph{Historical assessment}
	For each stock and year in the backtest window, the assessment will be re-run using only the information that would have been available at that time, including catches, discards, survey indices, biological data, and reference points. The base case will apply a standard SAM-style configuration across years that attempts to reconstruct the historical assessment as closely as is practically possible.
	
	\paragraph{Advice-rule}
	For each assessment, the advice rule in force in that year will be applied to generate catch advice for the following year. For the Celtic and Irish Sea gadoid stocks, this will be the standard ICES category~1 MSY rule, with fishing mortality set at \(F_{\mathrm{MSY}}\) above \(MSYB_{\mathrm{trigger}}\), reduced linearly as biomass declines towards \(B_{\lim}\), and set to zero where biomass is at or below \(B_{\lim}\). For Northeast Atlantic mackerel, the implemented rule will reflect the ICES advice, including any science-based management arrangements or stability constraints that applied in the relevant year.
	
	\paragraph{Counterfactual}
	Projections under perfect management for each annual assessment will be used as the starting point for short-term projections in which catches are aligned with the advised catches. This creates historical counterfactual trajectories that show what would have happened had scientific advice been followed without implementation error, political override, or TAC overshoot. Projections will track spawning-stock biomass, fishing mortality, catches, and the frequency of years spent in breach of precautionary thresholds and limits, or failing to achieve targets.
	
	\paragraph{Evaluation}
	Performance will be evaluated based on comparing realised and counterfactual performance (i.e. histories and trajectories) using a common set of metrics, e.g.
	
	\begin{itemize}
		\item \textbf{Biological status} including time spent below \(B_{\lim}\) and \(MSYB_{\mathrm{trigger}}\).
		\item \textbf{Exploitation} realised and counterfactual \(F\) relative to \(F_{\mathrm{MSY}}\).
		\item \textbf{Yield and stability} annual, cumulative catches and interannual variability.
	\end{itemize}
	
	\section*{Deliverables}	
	Deliverables will include:
	
	\begin{itemize}
		\item A week~1 screening summary, benchmark projections at target $F$, and confirmed list of case-study stocks.
		\item A technical report describing data sources, assessment reconstruction, advice-rule coding, projection methods, assumptions, diagnostics, and results by stock.
		\item Stock-specific summaries comparing realised and perfect-management trajectories of spawning-stock biomass, fishing mortality, catch, and advice.
		\item A cross-stock synthesis identifying where adherence to advice would have materially improved outcomes, where it would have had limited effect, and where shortcomings reflect the advice rule itself rather than implementation failure.
		\item Reproducible code and documentation to allow future updating of the backtest as new assessments become available.
	\end{itemize}
	
	\section*{Workplan}
	The study takes a step beyond standard retrospective analysis by separating two questions that are often conflated: whether the assessment and advice rule are robust, and whether management failed because advice was not followed. The results will support future decisions on harvest control rules and long-term management plans, including rebuilding strategies.
	
\textbf{Tasks (with indicative effort)}:
\begin{enumerate}
	\item \textbf{Initial screening and benchmark projections (5 days)}:
	obtain the most recent ICES stock assessments; extract historical stock
	estimates and advice from ICES Stock Assessment Graphs and evaluate
	advice--management mismatch; evaluate advice stability; confirm case-study
	stocks. Then run short-term projections at target $F$ from the most recent
	assessment and compare projected SSB, $F$, and catch outcomes with the
	current assessment as a benchmark; agree backtest period and deliverables.
	
	\item \textbf{Assessment reconstruction (20 days)}:
	re-run historical assessments and validate against published ICES outputs
	for the confirmed case-study stocks, including basic diagnostics.
	
	\item \textbf{Counterfactual simulation (15 days)}:
	apply advice rules and project forward under perfect implementation;
	generate counterfactual trajectories of SSB, $F$, and catches.
	
	\item \textbf{Comparative analysis and reporting (10 days)}:
	quantify differences between realised and counterfactual outcomes, prepare
	stock-specific summaries and cross-stock synthesis, and final report.
\end{enumerate}

	
\end{document}